\section{議論と未解決の問題}\label{sec:discussion}

\paragraph{全ての framing での正値性．}
個々の結び目では，$x$ の指数和とニュートン形によって，全ての $x$ での正値性を有限個の判定に帰着できた（定理~\ref{thm:two-all-x}）．
しかしこれは結び目ごとの判定で，一般の結び目についての定理ではない．
ニュートン形の係数の正値性を，入力の性質（ホモロジーの非負性と $\sigma$ 不変性など）から導けるかは未解決である．
ツイスト結び目で $p$ と $x$ が逆符号の領域と，3 箱の全ての $x$ の証明書も得られていない．
また，Cheng--Sułkowski~\cite{Cheng21} の toric brane の refined open BPS 不変量と，framing をつけた自明な結び目の $\hf_\rho$ の対応は確かめていない．

\paragraph{framing の物理的な導出．}
Mariño--Vafa の符号が Laurent 性に必要であることは示した（命題~\ref{prop:sign}）．
しかし，refined の平坦座標のずれ $V\mapsto-(q/t)^{1/2}V$ を幾何から導くことはしていない．
弦側で $U(1)_R$ の作用が framing によらず自明であることも仮定にとどまる．

\paragraph{双曲結び目の $[2,1]$．}
双曲結び目の 3 箱には，定義に基づく refined $[2,1]$ が必要である．候補は三つある．
(i) refined の Racah 行列による構成（arborescent 計算の refined 化）．
(ii) DAHA・楕円 Hall 代数型の定義の双曲結び目への拡張（種数 2 の DAHA は $\mathfrak{sl}_2$ に限られる~\cite{Arthamonov15}）．
(iii) $[2,1]$ 彩色 HOMFLY-PT ホモロジーの計算と，KN の $\rCS$ への変換規則．
いずれも本論文の範囲を超える．

\paragraph{正規化．}
単項式の正規化の閉じた式（計算結果~\ref{cr:M}）は，refined 衛星公式（仮説~\ref{hyp:sat}）の下で導出した．
DAHA と楕円 Hall 代数は単項式を決めないので，この仮説は物理的な仮説として残る．
トーラス結び目の部分では，Mellit の系の標準の正規化での形を計算で確かめた（定理~\ref{thm:M}）．
$d_\pm$ の表現の元と Bergeron--Garsia--Leven--Xin の $Q_{m,n}$ の対応を具体的に示すことが残っている．
ただし，これは正値性の結果には影響しない．

\paragraph{そのほか．}
pretzel 結び目の異符号の領域の残り（\S\ref{sec:two}）と，2 箱の単峰性の構造的な理由も未解決である．
